\section{Conclusion}\label{sec:conclusion}

Urban drone decisions depend on the city, the weather and the battery at once, and a decision that sees only part of
this picture fails in systematic, city- and weather-dependent ways: returns fly into buildings, modules contradict each
other in wind, and scans miss their targets without noticing. We presented Drone4D, a space-time world runtime that serves one model of
geometry, weather and energy to the planner, the fleet runtime and the browser, and a decision layer that couples the
launch precheck, return planning, scan planning and task delegation to it. Batched queries on one clock and a rotating
refresh make the coupling cheap enough for 1,000 drones on one CPU core. In closed-loop campaigns over seven cities and
twelve weather conditions, coupling reduced unsafe sorties from \PXUNSAFE\% to \CPUNSAFE\% while doubling the share of
completed sorties, kept sorties out of a forecast weather front, restored recognition capture in fog and raised the
success of task delegation. Each coupling removes a different failure, and none of them alone delivers both safety and
productivity. 
The runtime, the experiment harness, the raw results and the figure scripts are released as open source,
so that calibrated vehicle, weather and perception models can be plugged into the same shared interfaces.
